% ===================================================================
%  TAVOLA N. 1 — La scuola. --- Il liceo. --- L'aula.
%
%  Ceci n'est PAS une transcription : c'est une traduction, faite en
%  2026, en italien d'aujourd'hui. D'ou trois differences avec
%  texto/io et texto/fr, et trois seulement :
%
%   * pas de \nl ni de \cc — la traduction ne suit aucune ligne
%     imprimee, et les lignes de ce fichier ne sont que du confort ;
%   * pas de \begin{VUpage} — elle n'a ni feuillet ni folio, et la
%     page de lecture ne lui en affiche pas ;
%   * le gras se pose sur le TERME seul, ponctuation dehors.
%
%  Tout le reste est identique, et doit l'etre : les cles « %%K » sont
%  celles de texto/io, macro pour macro pour l'apparat, et les renvois
%  (n) visent les MEMES objets de la planche.
%
%  LA SOURCE EST DE 1926, LA COLONNE EST DE 2026. Voir LISEZ-MOI § 5.
%  Pour l'italien la regle mord surtout sur trois points : pas de
%  passe simple narratif la ou l'usage parle emploie le passe compose,
%  pas de « egli » ni de « ella » pour « lui » et « lei », et pas de
%  subjonctif de politesse la ou personne n'en met plus.
%
%  UNE DECISION DE COMPOSITION. L'italien ecrit « TAVOLA N. 1 ». Le
%  mot ne se confond avec aucun autre membre de NUMERO_TAB : le
%  « TABELLA » de l'interlingua et du romanche porte un B la ou
%  celui-ci porte un V, et le « TAULA » du catalan, de l'occitan et du
%  basque n'a pas de V du tout. Le titre du volume, « Tavole
%  ausiliarie Delmas », est en bas de casse. Le cablage a recu
%  « TAVOLA » et trois ordinaux ; « terza » y etait deja par le
%  romanche, « serie » et « scena » par le groupe commun.
%
%  ET CETTE COLONNE-CI EST UN MIROIR, CE QU'AUCUNE DES DIX-HUIT
%  PRECEDENTES N'ETAIT. L'ido n'est pas une langue trouvee : c'est une
%  langue faite, en 1907, a partir de six autres — le francais,
%  l'italien, l'espagnol, l'anglais, l'allemand et le russe. L'italien
%  est donc l'un des six parents, et la colonne de gauche lui renvoie
%  sans cesse ses propres mots, a peine deguises : tablo, benko,
%  libro, papero, inko, klaso, lampo, muziko, geometrio, sfero. La
%  colonne allemande avait montre que l'ido COMPOSE comme l'allemand ;
%  celle-ci montrera qu'il NOMME comme l'italien.
%
%  CE QUE LA PAGE MONTRE : DEUX OBJETS QUI PORTENT ENCORE LE NOM DE
%  LEUR ORIGINE, ET PERSONNE D'AUTRE QUE L'ITALIEN NE L'ENTEND.
%
%  1. LA CATTEDRA DE L'ALINEA 1 (6). C'est le meme mot que la
%     « cathedra » de l'eveque, le siege d'ou l'on enseigne, et que la
%     CATHEDRALE, qui n'est rien d'autre que l'eglise ou ce siege se
%     trouve. Le pupitre du professeur s'appelle donc, en italien
%     d'aujourd'hui, un trone episcopal. L'allemand dit « Pult », le
%     planche ; l'anglais « desk », de « discus », le plateau ;
%     l'estonien a emprunte « kateeder » sans plus rien entendre. Seul
%     l'italien sait encore de quoi il parle.
%
%  2. LA LAVAGNA DE L'ALINEA 2 (1). Le tableau noir s'appelle du nom
%     d'un village de Ligurie, Lavagna, dont les carrieres d'ardoise
%     fournissaient l'Europe. L'OBJET PORTE LE NOM DU LIEU D'OU LA
%     PIERRE VENAIT. C'est la regle du tableau 11 allemand — la langue
%     d'une institution est celle du pays qui l'a batie — appliquee
%     non a une institution mais a une MATIERE, et c'est la premiere
%     fois que la serie la rencontre sous cette forme.
%
%  ET UNE AMBIGUITE QUI N'EN EST PAS UNE. L'italien appelle « riga »
%  la regle de l'alinea 14 (50) et « riga » aussi la ligne ecrite de
%  l'alinea 9. Ce n'est pas une confusion : la regle est l'instrument
%  qui fait la ligne, et le nom est passe de l'outil au resultat. La
%  page emploie les deux a huit alineas de distance.
%
%  LES NOMS DE PERSONNE SONT ITALIENS : Enrico, Stefano, Carlo,
%  Ludovico, Antonio, Paolo, Giorgio, Edmondo, Emilio, Augusto,
%  Bernardo, Alberto, Giacomo, Cristiano, Ernesto, Francesco,
%  Maurizio, Giovanni, Pietro. Et la note (*) commence justement par
%  « Giovanni » : de la liste de sept formes que Guignon donne pour
%  montrer le desordre des langues, la premiere est italienne.
%
%  « postigo » (t01, renvoi 78) EST UN VASISTAS, ET LA PLANCHE LE DIT.
%  Le mot a ete verifie sur la gravure en 2026, non sur les
%  dictionnaires : la verriere de la classe a deux rangees de carreaux,
%  et dans la rangee HAUTE deux petits vantaux sont graves battants,
%  gonds sur un montant vertical, ouverts vers l'interieur. C'est le
%  seul ouvrant de la verriere, et il ne sert qu'a aerer.
%
%  D'OU CE QU'IL N'EST PAS, et ou seize colonnes s'etaient trompees :
%  ni une lucarne de toit (it. « lucernari »), ni un volet (cs, lt, lb),
%  ni un meneau ni un montant de chassis (ar, ca, es, oc), ni une
%  imposte fixe, ni le battant de la grande fenetre (bn), ni le
%  « ranma » de la cloison japonaise. Et « framuga » russe nomme
%  l'imposte entiere, quand le francais nomme le VANTAIL : c'est
%  « fortochka » qu'il fallait, comme l'ukrainien avait deja
%  « kvatyrka ».
% ===================================================================

%%K t01-apar-0 apar
\VUsaut{2.0mm}
\VUcentre{12.6pt}{120}{LIBRETTO ESPLICATIVO}
\VUsaut{3.2mm}
\VUcentre{8.4pt}{160}{PER LE}
\VUsaut{2.6mm}
\VUtitre{15.6pt}{0}{Tavole ausiliarie Delmas}
\VUsaut{3.4mm}
\VUornamento{9pt}
\VUsaut{5.0mm}
\VUcentre{13.2pt}{90}{PRIMA SERIE}
\VUsaut{3.0mm}
\VUfilet{16mm}
\VUsaut{5.6mm}

%%K t01-apar-1 apar
\VUtitre{15.0pt}{60}{TAVOLA N. 1}
\VUsaut{1.4mm}
\VUtitre{11.4pt}{0}{La scuola. --- Il liceo. --- L'aula.}
\VUsaut{2.2mm}
\VUfilet{20mm}
\VUsaut{6.4mm}

%%K t01-tit-1 sub
\VUpk{10.2pt}{40}{Descrizione generale.}
\VUsaut{3.0mm}

%%K t01-01-1 p
1. --- Questa tavola mostra un'\VUgras{aula} di \VUgras{liceo}. In quest'aula ci sono
otto \VUgras{tavoli} \textsuperscript{(18)} e otto \VUgras{banchi}
\textsuperscript{(32)}. Su ogni banco siedono tre alunni. Il \VUgras{professore}
\textsuperscript{(7)} è seduto su una \VUgras{sedia} \textsuperscript{(8)} dietro la
sua \VUgras{cattedra} \textsuperscript{(6)}.

%%K t01-02-1 p
2. --- A sinistra della cattedra c'è un \VUgras{armadio} a vetri
\textsuperscript{(15)}. A destra c'è la \VUgras{lavagna} \textsuperscript{(1)} con il
\VUgras{cancellino} \textsuperscript{(2)} e il \VUgras{gesso} \textsuperscript{(3)}.

%%K t01-02-2 p
Sopra la cattedra sono appese le \VUgras{tavole murali} \textsuperscript{(9, 11, 12)}
e una \VUgras{carta geografica} \textsuperscript{(10)}.

%%K t01-03-1 p
3. --- Non tutti gli alunni sono al loro \VUgras{posto} \textsuperscript{(17)}.
Giovanni \textsuperscript{(4)} è in piedi alla lavagna; tiene il cancellino e
cancella le \VUgras{figure geometriche}.

%%K t01-03-2 p
Pietro è vicino alla cattedra; raccoglie i \VUgras{compiti scritti}
\textsuperscript{(5)} e li porta al professore.

%%K t01-04-1 p
4. --- Questo professore insegna \VUgras{lingue moderne}. Insegna agli alunni a
parlare, leggere e scrivere le lingue straniere (\VUgras{tedesco}, \VUgras{inglese},
\VUgras{spagnolo}, \VUgras{italiano}, \VUgras{russo} o \VUgras{francese}).

%%K t01-05-1 p
5. --- L'aula è molto ampia e molto luminosa. Nella parete di fondo ci sono una larga
\VUgras{finestra}, due \VUgras{vasistas} \textsuperscript{(78)} e una
\VUgras{porta a vetri} per arieggiarla e illuminarla.

%%K t01-05-2 p
Quest'aula non è fredda. In mezzo, per riscaldarla, c'è un'ottima \VUgras{stufa}
\textsuperscript{(46)} con il suo \VUgras{tubo} \textsuperscript{(45)}.

%%K t01-05-3 p
Alla parete sono fissati gli \VUgras{attaccapanni} \textsuperscript{(58)}; a essi sono
appesi i berretti e i cappotti degli alunni.

%%K t01-tit-2 sub
\VUsaut{5.2mm}
\VUpk{10.2pt}{40}{Il cortile.}
\VUsaut{3.6mm}

%%K t01-06-1 p
6. --- L'\VUgras{orologio} \textsuperscript{(63)} segna le nove e cinque.

%%K t01-06-2 p
La \VUgras{ricreazione} \textsuperscript{(76)} è appena finita e la lezione ricomincia.
La ricreazione si fa nel \VUgras{cortile} \textsuperscript{(75)}. Vediamo qualche
alunno che gioca. Un \VUgras{sorvegliante} \textsuperscript{(74)} li tiene d'occhio.

%%K t01-07-1 p
7. --- Nel cortile vediamo anche altre persone: l'\VUgras{ispettore}
\textsuperscript{(73)}, il \VUgras{preside} \textsuperscript{(71)} e il
\VUgras{vicepreside} \textsuperscript{(72)}.

%%K t01-07-2 p
L'ispettore controlla le scuole, il preside dirige il liceo, il vicepreside
sorveglia le lezioni.

%%K t01-07-3 p
La quarta persona è il \VUgras{portinaio} \textsuperscript{(77)}. Porta sotto il
braccio un registro per segnare gli assenti.

%%K t01-08-1 p
8. --- In fondo al cortile ci sono gli \VUgras{attrezzi ginnici}
\textsuperscript{(70)} e il laboratorio di \VUgras{lavori manuali}
\textsuperscript{(69)}.

%%K t01-08-2 p
A destra della tavola cominciamo a vedere le aule di \VUgras{musica} e di
\VUgras{disegno}.

%%K t01-noto-1 noto
\VUnotes{143.2mm}{%
(*) Per i \textit{nomi di battesimo} si consiglia di scriverli anche nella forma
latina. È l'unico modo di evitare innumerevoli varianti. Altrimenti come scegliere
fra \textit{Giovanni, Jean,} \textit{Johann, Jan, Hans, John, Ivan} e così via?
Faremo forse un dizionario apposta per i nomi di battesimo? Prendiamo dal latino il
loro nominativo e diciamo: \VUgras{Ioannes, Ioanna; Iulius, Iulia; Iakobus;
Andreas; Lukas.} (Grammatica completa).%
}

%%K t01-tit-3 sub
\VUpk{10.2pt}{40}{Descrizione particolareggiata.}
\VUsaut{2.4mm}

%%K t01-tit-4 sub
\VUpk{10.2pt}{40}{I bravi alunni.}
\VUsaut{3.0mm}

%%K t01-09-1 p
9. --- Nel primo banco a destra della cattedra siedono \VUgras{Enrico} (*),
\VUgras{Stefano} e \VUgras{Carlo}.

%%K t01-09-2 p
Enrico è molto attento e diligente; ascolta il professore. Sul suo tavolo ha un
\VUgras{quaderno} \textsuperscript{(28)}, un \VUgras{libro} \textsuperscript{(29)},
un \VUgras{dizionario} \textsuperscript{(30)} e un \VUgras{astuccio}
\textsuperscript{(31)}. Il suo libro ha molte \VUgras{pagine}; ogni capitolo contiene
diversi \VUgras{capoversi} e ogni \VUgras{riga} molte \VUgras{parole}.

%%K t01-09-4 p
Il suo vicino Stefano \textsuperscript{(24)} è volenteroso. Si segna sul quaderno la
lezione per la prossima volta. Ha una bella \VUgras{calligrafia}. Il suo libro è
chiuso. Ne vediamo solo la \VUgras{copertina} \textsuperscript{(27)}. Accanto a lui
c'è il suo \VUgras{compasso} \textsuperscript{(26)}.

%%K t01-09-5 p
Carlo \textsuperscript{(23)} tiene il libro in mano; lo apre per ripassare la lezione.
Come ogni alunno, ha davanti a sé un \VUgras{calamaio} \textsuperscript{(25)}. È
ubbidiente.

%%K t01-10-1 p
10. --- Al tavolo accanto siedono \VUgras{Ludovico}, \VUgras{Antonio} e
\VUgras{Paolo}. Ludovico recita la sua \VUgras{lezione} \textsuperscript{(22)} e
risponde alle \VUgras{domande} del professore. Antonio \textsuperscript{(21)} è molto
distratto; a volte il professore lo punisce perfino. Paolo \textsuperscript{(20)} è un
buon \VUgras{compagno}, gentile e sempre pronto ad aiutare. È molto attento.

%%K t01-tit-5 sub
\VUsaut{4.2mm}
\VUpk{10.2pt}{40}{I cattivi alunni.}
\VUsaut{3.2mm}

%%K t01-11-1 p
11. --- Dietro Enrico siede \VUgras{Giorgio} \textsuperscript{(38)}. Non è un ragazzo
cattivo, ma è \VUgras{chiacchierone}, distratto e irrequieto. La sua
\VUgras{leggerezza} e la sua \VUgras{disobbedienza} creano \VUgras{disordine} durante
la lezione, e spesso si merita una severa \VUgras{ammonizione}. La sua
\VUgras{brutta copia} \textsuperscript{(35)} è sporca, ci fa sopra
\VUgras{macchie d'inchiostro} con il suo \VUgras{pennino} \textsuperscript{(34)}, e
per questo usa continuamente la sua \VUgras{gomma} \textsuperscript{(36)}; la sua
\VUgras{cartella} \textsuperscript{(33)} è strappata, perché è molto trascurato.
Perché porta il suo \VUgras{portapenne} \textsuperscript{(37)} alla lezione di lingue
moderne?

%%K t01-11-2 p
Il suo vicino Edmondo \textsuperscript{(39)} non lo sopporta. È vero che è un po'
pigro, ma sta zitto e fermo. Non chiacchiera; solo che, invece di ascoltare,
preferisce leggere i \VUgras{racconti} del suo \VUgras{libro di lettura}.

%%K t01-11-3 p
Il terzo alunno di questo banco è \VUgras{Ludovico} \textsuperscript{(40)}. È
volenteroso. Scrive su un \VUgras{foglio di carta} bianco. Davanti a sé ha messo la
sua \VUgras{carta assorbente} \textsuperscript{(41)}.

%%K t01-12-1 p
12. --- Gli alunni del secondo banco, a sinistra del professore, sono ancora molto
piccoli. Il primo, il piccolo \VUgras{Emilio}, non sa ancora scrivere bene; porta con
sé la sua \VUgras{lavagnetta} \textsuperscript{(42)}, il suo \VUgras{gessetto} e la
sua \VUgras{spugna} \textsuperscript{(43)}.

%%K t01-12-2 p
Accanto a lui siedono \VUgras{Augusto} e \VUgras{Bernardo} \textsuperscript{(44)}.
Quest'ultimo lavora bene, ma il suo \VUgras{vestiario} è molto trascurato.

%%K t01-13-1 p
13. --- Dietro di loro siedono tre \VUgras{pigri}. \VUgras{Alberto} non studia le sue
lezioni. \VUgras{Giacomo} \textsuperscript{(47)} dorme invece di ascoltare. La sua
\VUgras{pigrizia} gli procura \VUgras{rimproveri} e perfino \VUgras{punizioni}. Infine
\VUgras{Cristiano} \textsuperscript{(48)} è troppo sventato. Si \VUgras{comporta} male
e non fa nessuno sforzo per correggersi.

%%K t01-14-1 p
14. --- I suoi vicini, al contrario, si comportano bene. \VUgras{Ernesto}
\textsuperscript{(51)} ama l'\VUgras{ordine} e la \VUgras{pulizia}; usa sempre la sua
\VUgras{riga} \textsuperscript{(50)} e la sua \VUgras{matita} \textsuperscript{(49)}.
È un \VUgras{alunno esterno}.

%%K t01-14-2 p
\VUgras{Francesco} \textsuperscript{(52)} è \VUgras{convittore}; porta la
\VUgras{divisa} del liceo; lì mangia e lì dorme. È un buon \VUgras{compagno}.

%%K t01-14-3 p
Maurizio fa la punta alla sua \VUgras{matita} con il suo \VUgras{temperino}
\textsuperscript{(56)}; è molto accurato.

%%K t01-14-4 p
Porta con sé tutti i suoi libri, la sua \VUgras{grammatica} \textsuperscript{(55)}, il
suo \VUgras{taccuino} \textsuperscript{(54)} e perfino un \VUgras{sottomano}
\textsuperscript{(53)}. Il suo \VUgras{zaino} \textsuperscript{(57)} è per terra
dietro di lui.

%%K t01-14-5 p
I due tavoli in fondo all'aula sono occupati dai \VUgras{bravi alunni}
\textsuperscript{(19)}.

%%K t01-tit-6 sub
\VUsaut{4.6mm}
\VUpk{10.2pt}{40}{Disegno e musica.}
\VUsaut{3.4mm}

%%K t01-15-1 p
15. --- Nell'\VUgras{aula di disegno} ci sono bei \VUgras{modelli} appesi al muro e
una \VUgras{lampada} \textsuperscript{(62)} con un \VUgras{paralume}, appesa al
soffitto. Il \VUgras{maestro} \textsuperscript{(60)} è molto paziente; corregge i
\VUgras{disegni} \textsuperscript{(61)} degli alunni e insegna loro a disegnare un
\VUgras{busto} \textsuperscript{(59)} dal vero.

%%K t01-16-1 p
16. --- L'\VUgras{aula di musica} sembra molto interessante. Lì si impara a leggere
la \VUgras{musica}, a solfeggiare le \VUgras{note della scala}
\textsuperscript{(67)} scritte sul \VUgras{pentagramma} (do, re, mi, fa, sol, la,
si\,=\,C. D. E. F. G. A. B.), a riconoscere le \VUgras{chiavi} e a cantare belle
\VUgras{melodie}. Gli spartiti si mettono sul \VUgras{leggio}
\textsuperscript{(65)}. Uno degli alunni \VUgras{suona il pianoforte}
\textsuperscript{(64)} e il \VUgras{maestro di musica} \textsuperscript{(66)}
\VUgras{batte il tempo} \textsuperscript{(68)}.

%%K t01-17-1 p
17. --- Cominciamo a vedere un angolo del laboratorio di \VUgras{lavori manuali}
\textsuperscript{(69)}, dove i ragazzi possono imparare a piallare il legno e a
limare il ferro. Non tutti i licei ne hanno uno.

%%K t01-tit-7 sub
\VUsaut{5.0mm}
\VUpk{10.2pt}{40}{L'aula di scienze.}
\VUsaut{3.6mm}

%%K t01-18-1 p
18. --- Il professore di matematica insegna agli alunni la \VUgras{geometria},
l'\VUgras{aritmetica}, il \VUgras{calcolo} e il \VUgras{sistema metrico}. Sulla
tavola vediamo le principali figure geometriche: il \VUgras{cerchio}
\textsuperscript{(a)}, il \VUgras{quadrato} \textsuperscript{(b)}, il
\VUgras{triangolo} \textsuperscript{(c)}, la \VUgras{linea} \textsuperscript{(d)}.

%%K t01-18-2 p
Vediamo anche un'\VUgras{operazione}; non è un'\VUgras{addizione}, né una
\VUgras{sottrazione}, né una \VUgras{moltiplicazione}: è una \VUgras{divisione}
\textsuperscript{(e)}.

%%K t01-19-1 p
19. --- Sulla tavola del sistema metrico distinguiamo: il \VUgras{metro}
\textsuperscript{(a)}, la \VUgras{bilancia} \textsuperscript{(b)} con i suoi
\VUgras{pesi}, il \VUgras{litro} \textsuperscript{(e)}, il \VUgras{decalitro}
\textsuperscript{(f)}, il \VUgras{decilitro} e il \VUgras{metro cubo}
\textsuperscript{(d)}.

%%K t01-19-2 p
Gli alunni studiano anche le \VUgras{scienze naturali} \textsuperscript{(11)} ---
zoologia \textsuperscript{(a)}, botanica \textsuperscript{(b)}, geologia
\textsuperscript{(c)} --- e la \VUgras{storia} \textsuperscript{(12)}.

%%K t01-19-3 p
Nell'armadio ci sono gli strumenti di \VUgras{fisica} \textsuperscript{(15)} e di
\VUgras{chimica} \textsuperscript{(16)}.

%%K t01-19-4 p
Sopra l'armadio ci sono: una \VUgras{sfera} \textsuperscript{(14)}, un \VUgras{cono} e
un \VUgras{mappamondo} \textsuperscript{(13)}.

%%K t01-19-5 p
Sopra le tavole è appesa una \VUgras{carta} che rappresenta le cinque parti del
mondo: l'\VUgras{Europa} \textsuperscript{(g)}, l'\VUgras{Asia}
\textsuperscript{(e)}, l'\VUgras{Africa} \textsuperscript{(f)}, l'\VUgras{America}
\textsuperscript{(ab)} e l'\VUgras{Oceania} \textsuperscript{(h)}, e i due grandi
oceani, il \VUgras{Pacifico} \textsuperscript{(c)} e l'\VUgras{Atlantico}
\textsuperscript{(d)}.
\VUsaut{6.0mm}
\VUfilet{22mm}
